\clearpage\fancyhead[L]{\small PALLAS.}\SourceLeaf{354}
\begin{center}\textbf{PALLAS.}\par\vspace{4pt}\hrule\vspace{4pt}\fontsize{10}{12}\selectfont Monatliche Correspondenz. Vol. XXV, p. 389, April 1812.\par\vspace{4pt}\hrule\vspace{4pt}\end{center}

\begin{flushright}Göttingen, 15 April 1812.\end{flushright}
\par\indent I was very impatient for Pallas to reappear this year, to see whether the perturbation calculations would be borne out as well this time as they were last year. But moonlight and bad weather kept me from seeing it until 4 April. On that evening the sky was clear for only a short time, and I only caught sight of it. It was not until 9 April that I managed to make an observation, and even then the declination was only estimated.

\begin{center}\fontsize{10}{12}\selectfont\setlength{\tabcolsep}{0pt}\renewcommand{\arraystretch}{1.15}\begin{tabular}{l}
1812 April 9\quad $11^{\mathrm u}\;32^{\mathrm m}\;27^{\mathrm s}$ R.A. $\PallasSymbol\;269^\circ\;6'\;6''7$.\quad Decl. $16^\circ\;38'$
\end{tabular}\end{center}

\par\indent This observation at least shows that the ephemeris calculated by Mr. \textsc{Nicolai} is entirely accurate. It is much to be hoped that the opposition on 10 June will not be missed. But observing with fixed instruments will not be easy because the planet is so faint. At present it is only of the 11th magnitude.

\par\indent I have now completely finished the first calculation of the periodic and secular perturbations caused by Jupiter. The total number of equations worked out is close to 400, and the largest, which depends on $5\,\JupiterSymbol-2\,\PallasSymbol$, amounts to about 56 minutes. It takes several days’ work to calculate a single position of Pallas. Some results of this work are of great interest.

\par\indent The sky was overcast, so nothing could be seen here of yesterday’s occultation of $\alpha$ Tauri.
